% Generated draft: requires mathematical and source review.
\documentclass[11pt]{article}
\usepackage[margin=25mm]{geometry}
\usepackage{fontspec}
\setmainfont{Latin Modern Roman}
\usepackage{amsmath,amssymb,mathtools,unicode-math}
\setmathfont{Latin Modern Math}
\usepackage{enumitem,needspace,xurl,hyperref}
\setlength{\parindent}{0pt}
\setlength{\parskip}{5pt}
\setlength{\emergencystretch}{3em}
\newcommand{\N}{\mathbb{N}}
\newcommand{\Z}{\mathbb{Z}}
\newcommand{\Q}{\mathbb{Q}}
\newcommand{\R}{\mathbb{R}}
\newcommand{\C}{\mathbb{C}}
\newcommand{\F}{\mathbb{F}}
\newcommand{\A}{\mathbb{A}}
\newcommand{\PP}{\mathbb{P}}
\newcommand{\E}{\mathbb{E}}
\newcommand{\eps}{\varepsilon}
\newcommand{\dd}{\,\mathrm{d}}
\begin{document}
\section*{UnsolvedMath ID: 30003048}
\textbf{Source problem code:} catalogue entry
\subsection{Definitions and mathematical statement}
Write \(p\colon M\to B=M/G\) for the principal \(G\)-bundle appearing in the source. Thus \(G\) acts freely on \(M\), the fibers of \(p\) are the \(G\)-orbits, and the bundle is locally trivial. The source does not specify the category of bundles, the class of groups \(G\), or compactness assumptions.

In a smooth geometric interpretation, a Dirac-type operator acts on sections of a Clifford module over a Riemannian manifold and is a first-order differential operator with principal symbol given by Clifford multiplication, subject to the chosen symbol convention. The source does not specify the metrics, Clifford modules, connections, or operator domains needed for such an interpretation. For a Hilbert-space realization \(D\) and a unitary implementation \(U_g\) of the group action, equivariance means
\[
U_g\operatorname{Dom}(D)=\operatorname{Dom}(D),\qquad DU_g\xi=U_gD\xi
\quad(g\in G,\ \xi\in\operatorname{Dom}(D)).
\]
The required lift of the action to the geometric operator data is not supplied.

For terminology, a spectral triple \((\mathcal A,H,D)\) consists of a represented \(*\)-algebra \(\mathcal A\) on a Hilbert space \(H\) and a densely defined self-adjoint operator \(D\), with \(a\operatorname{Dom}(D)\subseteq\operatorname{Dom}(D)\), bounded extensions of \([D,a]\), and compact operators \(a(1+D^2)^{-1/2}\) for \(a\in\mathcal A\). In the unital case with a unital representation, the latter condition gives compact resolvent. A usual equivariant version additionally has compatible actions \(\alpha_g\) on \(\mathcal A\) and \(U_g\) on \(H\), satisfying \(U_gaU_g^{-1}=\alpha_g(a)\) and the equivariance conditions above. These conventions provide terminology only; the source does not select a precise class of equivariant spectral triples.

\textbf{Research question.} Formulate a mathematically precise meaning of factorization of an equivariant Dirac operator on \(p\colon M\to M/G\) into Dirac operators associated with \(G\) and \(M/G\), and formulate a generalization of that meaning to noncommutative equivariant spectral triples.

The requested notion of factorization is itself unknown in this formulation. The source specifies neither an operator identity nor a tensor-product construction, an equivalence relation, allowable correction terms, or hypotheses guaranteeing the existence of the factor operators. Consequently, this is a request for a framework, not a fully quantified assertion that every principal bundle or every equivariant spectral triple admits a particular factorization.
\subsection{Short English statement}
The question asks how to express a Dirac operator on a principal \(G\)-bundle in terms of operators associated with the group fibers and the quotient, and how to extend this idea to noncommutative equivariant spectral triples. The meaning of factorization and its necessary hypotheses remain part of the question.
\subsection{Sources}
\begin{itemize}
\item Supplied local UnsolvedMath registry record, permanent ID \(30003048\), exact source problem code \texttt{catalogue entry}; title: \emph{Factorization of Equivariant Dirac Operators and Spectral Triples}. Collection: Oberwolfach Reports Open Problems. The record gives the proposed year as 2016 and supplies no individual proposer.
\item \emph{Mini-Workshop: Operator Spaces and Noncommutative Geometry in Interaction} (2016), Oberwolfach Reports, pp.\ 267--295, DOI 10.4171/owr/2016/6. Supplied source citation and literature-source link, labeled ``original OWR report'': \url{https://doi.org/10.4171/owr/2016/6}. The background describes the report as documenting workshop discussions and open problems.
\item Literature assessment dated \texttt{2026-08-21}: the supplied curation labels the problem open and characterizes it as a research direction concerning a general meaning and framework for Dirac-operator factorization, with no general resolution provided by the cited discussion. This is supplied dated curation, not an independently verified literature assessment; the record explicitly requests independent verification of time-sensitive status and cited claims.
\item The supplied statement-status label is \texttt{corrected\_verified}. Its accompanying verification note reports only safe notation, typography, and TeX-delimiter cleanup, with hypotheses, quantifiers, and subquestions preserved. This label is not independent mathematical verification of the present draft.
\end{itemize}
\end{document}
